\subsubsection{Phân tích VC Dimension và Structural Risk Minimization}

\paragraph{VC Dimension của bộ phân lớp tuyến tính}

Xét lớp giả thuyết $\mathcal{H}$ gồm các bộ phân lớp tuyến tính trong không gian $\mathbb{R}^D$:
\[
f(\boldsymbol{x}) = \text{sign}(\boldsymbol{w}^T \boldsymbol{x} + b)
\]

VC dimension (Vapnik–Chervonenkis dimension) của $\mathcal{H}$ là số lượng điểm lớn nhất có thể được \textit{shatter} bởi lớp giả thuyết này. Một tập điểm được gọi là shatter nếu với mọi cách gán nhãn nhị phân, luôn tồn tại một siêu phẳng phân tách đúng.

Một kết quả kinh điển trong lý thuyết học máy cho thấy:
\[
\text{VC}(\mathcal{H}) = D + 1
\]

Điều này có nghĩa là:
\begin{itemize}
    \item Với $D+1$ điểm tổng quát, luôn tồn tại một siêu phẳng phân tách mọi cách gán nhãn
    \item Nhưng với $D+2$ điểm, không phải mọi cấu hình nhãn đều có thể phân tách tuyến tính
\end{itemize}

Do đó, VC dimension phản ánh trực tiếp độ phức tạp (capacity) của mô hình tuyến tính.

\paragraph{Liên hệ với khả năng tổng quát hóa}

Theo lý thuyết học thống kê, VC dimension cho phép xây dựng các chặn trên cho sai số tổng quát hóa. Cụ thể, với xác suất cao, ta có:
\[
R(f) \le \hat{R}(f) + O\left(\sqrt{\frac{\text{VC}(\mathcal{H})}{N}}\right)
\]

trong đó:
\begin{itemize}
    \item $R(f)$ là rủi ro thực (true risk)
    \item $\hat{R}(f)$ là rủi ro thực nghiệm (empirical risk)
    \item $N$ là số mẫu
\end{itemize}

Hệ quả quan trọng:
\begin{itemize}
    \item VC dimension càng lớn $\Rightarrow$ mô hình càng linh hoạt nhưng dễ overfitting
    \item VC dimension nhỏ $\Rightarrow$ mô hình đơn giản nhưng có thể underfitting
\end{itemize}

Do đó, kiểm soát độ phức tạp mô hình là yếu tố then chốt để đảm bảo khả năng tổng quát hóa.

\paragraph{Structural Risk Minimization (SRM)}

Structural Risk Minimization (SRM) là một nguyên lý lựa chọn mô hình dựa trên việc cân bằng giữa:
\begin{itemize}
    \item Sai số huấn luyện (empirical risk)
    \item Độ phức tạp mô hình (model capacity)
\end{itemize}

Ý tưởng cốt lõi của SRM là xây dựng một dãy các lớp giả thuyết lồng nhau:
\[
\mathcal{H}_1 \subset \mathcal{H}_2 \subset \cdots \subset \mathcal{H}_m
\]

trong đó độ phức tạp (ví dụ VC dimension) tăng dần. Sau đó chọn lớp tối ưu bằng cách tối thiểu hóa chặn trên của rủi ro:
\[
\hat{R}(f) + \Omega(\mathcal{H})
\]

với $\Omega(\mathcal{H})$ là hàm phạt phụ thuộc vào độ phức tạp của lớp giả thuyết.

\paragraph{Liên hệ với Logistic Regression và tham số $C$}

Trong Logistic Regression với regularization $L_2$, tham số $C$ (nghịch đảo của $\lambda$) đóng vai trò điều khiển độ phức tạp hiệu dụng của mô hình:
\begin{itemize}
    \item $C$ nhỏ $\Rightarrow$ regularization mạnh $\Rightarrow$ mô hình đơn giản
    \item $C$ lớn $\Rightarrow$ regularization yếu $\Rightarrow$ mô hình phức tạp hơn
\end{itemize}

Mặc dù VC dimension lý thuyết của bộ phân lớp tuyến tính là $D+1$, nhưng regularization làm giảm \textit{effective capacity} của mô hình, từ đó cải thiện khả năng tổng quát hóa.

Do đó, việc thay đổi $C$ có thể được xem như xây dựng một chuỗi các mô hình với độ phức tạp tăng dần, phù hợp với nguyên lý SRM.

\paragraph{Kết quả thực nghiệm và phân tích SRM}

Thực nghiệm được tiến hành trên tập dữ liệu Covtype với:
\[
D = 54, \quad \text{VC dimension} = D + 1 = 55
\]

Mô hình sử dụng là Logistic Regression với regularization $L_2$, trong đó tham số $C = \frac{1}{\lambda}$ điều khiển độ phức tạp hiệu dụng của mô hình. Các giá trị $C$ được lựa chọn theo thang log:
\[
C \in \{10^{-3}, 10^{-2}, 10^{-1}, 1, 10, 10^2\}
\]

\paragraph{Kết quả}

Hiệu năng của mô hình được đo bằng accuracy trên tập train, validation và test:

\[
\begin{array}{c|c|c|c}
C & \text{Train} & \text{Validation} & \text{Test} \\
\hline
0.001 & 0.7174 & 0.7161 & 0.7199 \\
0.01  & 0.7234 & 0.7219 & 0.7257 \\
0.1   & 0.7247 & 0.7228 & 0.7264 \\
1     & 0.7246 & 0.7227 & 0.7265 \\
10    & 0.7246 & 0.7229 & 0.7262 \\
100   & 0.7246 & 0.7228 & 0.7264 \\
\end{array}
\]

Giá trị tối ưu theo validation:
\[
C^* = 10, \quad \text{Validation Accuracy} = 0.7229
\]

\paragraph{Biểu diễn trực quan}

\begin{figure}[htbp]
    \centering
    \includegraphics[width=0.8\linewidth]{srm_covtype.png}
    \caption{Hiệu năng của Logistic Regression theo tham số $C$ trên tập Covtype.}
    \label{fig:srm_covtype}
\end{figure}

\paragraph{Phân tích xu hướng}

Từ Bảng kết quả và Hình \ref{fig:srm_covtype}, có thể rút ra các nhận xét quan trọng:

\begin{itemize}
    \item Khi $C$ nhỏ ($10^{-3}$), regularization mạnh làm mô hình bị \textit{underfitting}, thể hiện qua accuracy thấp trên cả train và validation.

    \item Khi $C$ tăng từ $10^{-3}$ đến $10^{-1}$, accuracy tăng nhanh. Điều này cho thấy mô hình dần thoát khỏi trạng thái underfitting và học được cấu trúc dữ liệu hiệu quả hơn.

    \item Từ $C \ge 10^{-1}$, các đường cong trong Hình \ref{fig:srm_covtype} bắt đầu \textbf{bão hòa}. Cụ thể:
    \[
    \text{Validation Accuracy} \approx 0.723
    \]
    cho thấy mô hình đã đạt đến giới hạn biểu diễn trong không gian đặc trưng tuyến tính.

    \item Khoảng cách giữa đường train và validation rất nhỏ (dưới $0.002$), và đường test gần như trùng với validation, cho thấy:
    \begin{itemize}
        \item Không có hiện tượng overfitting đáng kể
        \item Mô hình tổng quát hóa tốt
    \end{itemize}

    \item Khi $C$ lớn ($\ge 10$), việc tăng độ phức tạp không mang lại cải thiện đáng kể, thể hiện qua việc các đường cong gần như phẳng.
\end{itemize}

\paragraph{Diễn giải dưới góc nhìn SRM}

Theo nguyên lý Structural Risk Minimization:
\[
R(f) \le \hat{R}(f) + \Omega(\mathcal{H})
\]

Trong đó:
\begin{itemize}
    \item $\hat{R}(f)$: sai số huấn luyện
    \item $\Omega(\mathcal{H})$: độ phức tạp mô hình
\end{itemize}

Kết quả thực nghiệm cho thấy:

\begin{itemize}
    \item Khi $C$ nhỏ: $\Omega$ nhỏ nhưng $\hat{R}$ lớn $\Rightarrow$ underfitting
    
    \item Khi $C$ tăng: $\hat{R}$ giảm nhanh trong khi $\Omega$ tăng vừa phải
    
    \item Khi $C$ lớn: cả $\hat{R}$ và validation error gần như không đổi $\Rightarrow$ đạt cân bằng tối ưu giữa bias và variance
\end{itemize}

Điều này phản ánh chính xác cơ chế của SRM trong việc lựa chọn mô hình tối ưu.

\paragraph{Vai trò của VC Dimension}

Mặc dù VC dimension lý thuyết của mô hình là:
\[
\text{VC} = 55
\]

trong thực tế, hiệu năng tổng quát hóa còn phụ thuộc mạnh vào kích thước dữ liệu. Với $N$ rất lớn, ta có:
\[
O\left(\sqrt{\frac{\text{VC}}{N}}\right) \ll 1
\]

Do đó:
\begin{itemize}
    \item Sai số tổng quát hóa được kiểm soát chặt chẽ
    \item Mô hình ít bị overfitting ngay cả khi giảm regularization
\end{itemize}

Ngoài ra, regularization đóng vai trò giảm \textit{effective capacity} của mô hình, mặc dù không làm thay đổi VC dimension lý thuyết.

\paragraph{Nhận xét quan trọng}

\begin{itemize}
    \item Ảnh hưởng của $C$ chỉ rõ rệt trong giai đoạn đầu, sau đó giảm dần.
    
    \item Điều này cho thấy giới hạn chính của mô hình không nằm ở regularization mà ở:
    \begin{itemize}
        \item Giả định tuyến tính
        \item Khả năng biểu diễn của đặc trưng
    \end{itemize}

    \item Việc tăng độ phức tạp mô hình (qua $C$) không thể vượt qua giới hạn này.
\end{itemize}

\paragraph{Kết luận}

\begin{itemize}
    \item Logistic Regression với regularization $L_2$ hoạt động ổn định trên dữ liệu lớn
    
    \item SRM giúp lựa chọn giá trị $C$ tối ưu một cách hiệu quả
    
    \item Tuy nhiên, khi dữ liệu đủ lớn, vai trò của regularization giảm đáng kể
    
    \item Để cải thiện hiệu năng hơn nữa, cần sử dụng:
    \begin{itemize}
        \item Mô hình phi tuyến
        \item Hoặc feature engineering
    \end{itemize}
\end{itemize}


